\bookchapter{第二十一章：移动网络与信息安全深探}

\chapterintro{咱们每天用手机上网、支付、聊天，感觉方便得很。但你有没有想过，你的聊天记录、支付密码、照片……这些信息在网上“裸奔”的时候多吗？说实话，大部分人根本没想过这个问题，但一旦出事——比如账号被盗、银行卡被刷——那时候再后悔就晚了。这一章咱们就好好聊聊，怎么让你的数字生活不那么“透明”。}

\booksection{一、VPN：网络里的“隐形斗篷”}

VPN，全称虚拟专用网络（Virtual Private Network）。你可以把它想象成在你家和你要去的网站之间，建起了一条专属的、加密的地下隧道。你在隧道里走来走去，外面的人看不见你手里拿着什么、要去哪。

\begin{itemize}
  \item \textbf{普通上网}：你的手机公共网络目标网站。中间经过的路由器、交换机，随便哪个环节都可能有人"偷看"你发送的数据。尤其是在咖啡厅、机场那种公共Wi-Fi上，说难听点，简直是"裸奔"。
  \item \textbf{使用VPN}：你的手机VPN加密隧道VPN服务器目标网站。数据被加密封装，"偷窥者"只能看到一堆乱码。另外，目标网站看到的是VPN服务器的IP地址，而不是你的真实IP。
\end{itemize}

那VPN到底能干嘛？简单说：第一，\textbf{保护你在公共Wi-Fi上的隐私}；第二，\textbf{绕过地域限制}（比如看某些只在特定国家上线的内容）；第三，\textbf{防止你的网络运营商（ISP）记录你的浏览习惯}。不过我得提醒一句：\textbf{不要用免费VPN！}免费的背后往往意味着你的数据就是他们的“商品”。我有个朋友用了某个免费VPN，结果手机莫名其妙开始弹广告，后来发现那个VPN把他的浏览数据打包卖给广告商了。所以选VPN，要么选付费的大品牌（如NordVPN、ExpressVPN），要么干脆别用。

另外，VPN不是万能的。它只能保护“传输中”的数据，如果目标网站本身不安全、或者你的手机上已经中了恶意软件，VPN也帮不了你。

\booksection{二、公共Wi-Fi的“甜蜜陷阱”——比你想象的危险}

你可能觉得，不就是连个免费Wi-Fi嘛，能有多大风险？我给你讲个真事：安全研究人员曾经做过实验，在机场摆了一个叫“Starbucks\_Free”的Wi-Fi热点，不到一小时就有上百人连接。而实际上这个热点是研究人员自己搭的，所有经过这个热点的\textbf{未加密流量}——比如你访问的HTTP网页内容、未启用TLS的App通信——他们确实能截获和查看。

但要注意一个关键前提：今天你日常用的微信、网银、支付宝，绝大多数都已经走了HTTPS加密。说白了，攻击者即使截获了数据包，看到的也只是一堆乱码，\textbf{无法直接读出你的聊天内容或支付密码}。不过他们依然能知道你访问了哪些网站（域名），如果某些老旧App或网页还在用HTTP明文传输，那内容就会完全暴露。这种攻击叫“中间人攻击”（Man-in-the-Middle），原理是攻击者在你和真正的Wi-Fi路由器之间“插了一脚”。更可怕的是，很多人的手机设置了“自动连接已知Wi-Fi”，攻击者只要把自己的热点名字设成你之前连过的某个Wi-Fi（比如“ChinaNet”或“CMCC”），你的手机就会自动连上去。

那怎么防？几条实用建议：

\begin{itemize}
  \item \textbf{认准HTTPS}：地址栏开头是https://，多出的"s"代表Secure。访问任何需要输密码的网站时，如果看到的是http://（没有s），立刻停下来别输！现在主流浏览器都会在非HTTPS网站前显示"不安全"的警告，看到这个标记就要警惕。
  \item \textbf{关掉"自动连接"}：不管是iPhone还是安卓，都在Wi-Fi设置里把"自动加入"或"自动连接"关掉。宁可每次手动点一下，也别让手机自己判断。
  \item \textbf{用手机流量做敏感操作}：转账、登录网银、输入身份证号……这些操作老老实实用自己的4G/5G流量，别省那点流量钱。
  \item \textbf{忘记不常用的Wi-Fi}：出差回来、旅游回来，把之前临时连过的酒店Wi-Fi、机场Wi-Fi都"忘记"掉。这样下次手机不会自动连到同名假热点。
\end{itemize}

\booksection{三、HTTPS是怎么保护你的？——SSL/TLS握手}

刚才提到了HTTPS，那它到底是怎么工作的呢？每次你访问一个HTTPS网站，浏览器和服务器之间会进行一次“SSL/TLS握手”，整个过程大概就像这样：

\begin{enumerate}
  \item 你的浏览器说："嘿，服务器，我想跟你安全通信，这是我的加密能力清单（支持的加密算法等）。"
  \item 服务器回复："好的，我从你的清单里选了这套加密方案，这是我的数字证书，里面包含我的公钥，由某某CA机构签名认证。"
  \item 浏览器验证证书的真假：是不是由可信的CA签发？有没有过期？域名对不对得上？
  \item 验证通过后，双方用公钥加密协商出一个"会话密钥"，之后的通信全用这个密钥加密。
\end{enumerate}

这就像你去银行办事：你先递上身份证（证书），柜员核实你的身份，然后你们俩约定一个暗号，之后的对话只有你俩能听懂。即使有人在旁边偷听，也只能听到加密后的乱码。

但注意：HTTPS只保证“传输过程”是加密的，不保证目标网站是好人。钓鱼网站也可以用HTTPS——那个小锁图标只能说明你连接的是“某个网站”，不代表那个网站就是你想去的那个。所以还是要自己核对网址是否正确。

\booksection{四、DNS安全：你的“网络电话本”安全吗？}

DNS（域名系统）就像是互联网的电话本：你输入“baidu.com”，DNS帮你查到它对应哪个IP地址（比如110.242.68.66），然后浏览器才能连过去。问题是，传统的DNS查询是\textbf{明文}的、\textbf{没有加密}的！等于说，你的ISP、Wi-Fi提供者、或者任何中间节点，都能看到你访问过哪些网站。

更恶心的是“DNS劫持”：攻击者在DNS环节做手脚，把“baidu.com”指向一个假的钓鱼网站。你明明输入的是正确网址，却被带到了一个一模一样但实际是假的登录页面，一输入账号密码就完蛋了。

好在现在有了加密DNS方案：

\begin{itemize}
  \item \textbf{DNS over HTTPS（DoH）}：把DNS查询藏在HTTPS流量里，混在正常的网页浏览流量中，根本分不出来哪些是DNS查询。
  \item \textbf{DNS over TLS（DoT）}：用TLS加密协议专门保护DNS查询，端口独立，网管能知道你用了DoT但看不到内容。
\end{itemize}

我个人推荐用Cloudflare的1.1.1.1或者Google的8.8.8.8作为DNS服务器，它们都支持加密DNS。在安卓的“私人DNS”设置里直接填“one.one.one.one”或者“dns.google”就行，超简单。iPhone稍微麻烦点，需要装个描述文件或者通过App设置，但也值得搞一下。

\booksection{五、“无痕浏览”真的能隐身吗？}

先说结论：\textbf{无痕模式是你家窗帘，不是隐身斗篷。}

它做了什么：不在本地留下痕迹——关闭无痕窗口后，浏览器不保存你的浏览历史、Cookies、表单填写记录、搜索记录。借别人电脑查个东西、或者买礼物不想被家人从浏览记录发现，用无痕模式是合适的。

它做不到什么：

\begin{itemize}
  \item 你访问的网站\textbf{照样能看到你来了}，你的IP地址、设备信息、浏览行为，网站端一清二楚。
  \item 你的网络运营商\textbf{照样能看到你访问了哪些网站}，只是看不到具体内容（如果用了HTTPS）。
  \item 公司/学校的网络管理员\textbf{一样能监控}你的上网行为。
  \item 你下载的文件、添加的书签会\textbf{保留下来}，不会随着关闭窗口而消失。
\end{itemize}

所以记住：无痕模式是“本地擦除器”，不是“网络隐身衣”。如果你真的需要隐藏上网痕迹，得用VPN+Tor浏览器这样的组合拳，不过对于日常使用来说，了解这些就够用了。

\booksection{六、浏览器隐私：别让浏览器把你卖了}

你可能没意识到，浏览器本身就是一个巨大的信息收集器。Cookies、浏览历史、缓存、自动填充、位置权限……每一样都可能泄露你的隐私。几个建议：

\begin{itemize}
  \item \textbf{定期清理浏览数据}：每隔一段时间，去设置里把Cookies、缓存、历史记录清一清。不是为了防止"被追踪"——那个靠清Cookies已经不太管用了——而是为了防止手机越来越卡。浏览器缓存能占好几个G的空间。
  \item \textbf{关掉"请勿追踪"的幻想}：浏览器里的"发送'请勿追踪'请求"选项，本质上是"请"网站不要追踪你。但网站听不听完全是另一回事，大部分网站直接忽略。别指望它。
  \item \textbf{善用隐私保护浏览器}：Firefox Focus、DuckDuckGo浏览器、Brave这些浏览器默认拦截追踪器和广告，比Chrome和Safari的默认设置要干净得多。
  \item \textbf{搜索引擎也重要}：用DuckDuckGo代替百度/Google搜索，它不记录你的搜索历史，不建立用户画像。说实话搜索结果确实不如Google精准，但日常查资料够用了。
  \item \textbf{检查App权限}：去手机设置里看看你的浏览器被授予了什么权限。它真的需要访问你的位置、麦克风、通讯录吗？大概率不需要。
\end{itemize}

\booksection{七、实用安全清单：保护自己从现在开始}

说了这么多，最后给你一个简单粗暴的安全清单，花半小时搞定，你的手机安全性会有质的飞跃：

\begin{enumerate}
  \item \textbf{关掉Wi-Fi自动连接}：iPhone在设置-无线局域网里，安卓在Wi-Fi设置里。
  \item \textbf{换一个加密DNS}：安卓直接设置私人DNS为"dns.google"，iPhone用App或描述文件设置1.1.1.1的DoH。
  \item \textbf{开启双因素认证（2FA）}：微信、支付宝、Apple ID、Google账号，能开2FA的都开。短信验证码已经是最低级别的2FA了，能用TOTP（如Google Authenticator、Microsoft Authenticator）更好。
  \item \textbf{查一遍App权限}：把每个App的权限过一遍，特别是位置、通讯录、短信、麦克风、相机这五项。你会发现有些手电筒App居然要你的通讯录——赶紧卸载。
  \item \textbf{定期清理浏览器数据}：设置个日历提醒，每月一次，清缓存、Cookies。
  \item \textbf{不要用公共充电桩的数据线}：USB口既能充电也能传数据，恶意充电桩（叫"Juice Jacking"）可以趁你充电时偷数据。出门带自己的充电头和线。
\end{enumerate}

\begin{tcbpractice}{现在就做：给你的手机做一次"安全体检"}
\begin{enumerate}
  \item \textbf{DNS升级}：安卓：设置 > 连接与共享 > 私人DNS > 填入"one.one.one.one"或"dns.google"。iPhone：用Safari搜索"1.1.1.1 DNS profile"安装描述文件，或下载Cloudflare的1.1.1.1 App。
  \item \textbf{Wi-Fi审计}：打开Wi-Fi列表，把不认识的、不常用的、名字怪怪的热点都"忘记"掉。
  \item \textbf{权限大扫除}：进入设置 > 隐私/权限管理，逐个App审查权限。标记出那些权限要求不合理（比如计算器要通讯录）的App，考虑卸载或用替代品。
  \item \textbf{检查账号安全}：打开Apple ID或Google账号的安全页面，看看有没有不明设备登录，确认2FA已开启，检查恢复联系人是否可信。
\end{enumerate}

\end{tcbpractice}
